\section{Introduction}\label{sec:intro}
A connection Laplacian couples the geometry of a manifold to parallel
transport in a vector bundle. Its boundary measurements therefore pose
two simultaneous inverse problems: determining the metric that governs
its principal part and determining the connection that governs the
transport of vector-valued solutions. We consider both unknowns at zero
frequency, with all measurements confined to one boundary patch.

Let $M$ be a compact connected smooth manifold of dimension $n\geq3$,
with nonempty smooth boundary, and let $\Gamma\subset\partial M$ be a
nonempty relatively open proper subset. The bundle is the trivial
Hermitian bundle $M\times\C^2$. A smooth unitary connection is written
$d_A=d+A$, where
\[
 A\in C^\infty(M;T^*M\otimes\mathfrak u(2)),\qquad
 \mathfrak u(2)=\{B\in\C^{2\times2}:B^*=-B\}.
\]
For a smooth Riemannian metric $g$, write
$L_{g,A}=d_A^{*g}d_A$. Its zero-Dirichlet realization is invertible:
a section of zero energy is parallel, and a parallel section with zero
boundary trace is zero. Elliptic coercivity then gives, for each
$f\in C_c^\infty(\Gamma;\C^2)$ extended by zero to $\partial M$, a unique
solution $u_f^{g,A}$ of
\[
 L_{g,A}u_f^{g,A}=0\quad\text{in }M,\qquad
 u_f^{g,A}|_{\partial M}=f.
\]
The measured object is the sesquilinear energy form
\begin{equation}\label{eq:DN}
 \langle\Lambda_{g,A,\Gamma}f,h\rangle
 =\int_M\langle d_Au_f^{g,A},d_Au_h^{g,A}\rangle_g\,dV_g,
 \qquad f,h\in C_c^\infty(\Gamma;\C^2).
\end{equation}
The Hermitian product is linear in its first argument. All integrations
use densities; no orientation of $M$ is required. In particular,
\eqref{eq:DN} specifies the boundary measurement without presupposing a
boundary metric or a boundary measure.

\begin{theorem}\label{thm:main}
Let $M$ and $\Gamma$ be as above. For $j=1,2$, let $g_j$ be a smooth
Riemannian metric on $M$ and let
$A_j\in C^\infty(M;T^*M\otimes\mathfrak u(2))$. If
\[
 \langle\Lambda_{g_1,A_1,\Gamma}f,h\rangle
 =\langle\Lambda_{g_2,A_2,\Gamma}f,h\rangle
 \quad\text{for all }f,h\in C_c^\infty(\Gamma;\C^2),
\]
then there are a smooth diffeomorphism $\Phi:M\longrightarrow M$ and
a smooth map $U:M\longrightarrow U(2)$ such that
\begin{equation}\label{eq:main-conclusion}
 \Phi|_\Gamma=\Id,\qquad U|_\Gamma=I,\qquad
 g_2=\Phi^*g_1,\qquad
 A_2=U^{-1}(\Phi^*A_1)U+U^{-1}dU.
\end{equation}
Here $I$ denotes the $2\times2$ identity matrix.
\end{theorem}
The two transformations in \eqref{eq:main-conclusion} preserve
\eqref{eq:DN}: the corresponding change of solution is
$u_2=U^{-1}(u_1\circ\Phi)$. Thus the conclusion describes precisely the
ambiguities inherent in these measurements. The patch need not be
connected, and the manifold may have several boundary components.

\subsection{Context and contribution}
Calder\'on's inverse conductivity problem asks whether electrical
measurements at the boundary determine an interior conductivity
\cite{Calderon1980}. For smooth isotropic conductivities in dimension
at least three, global uniqueness was established by Sylvester and
Uhlmann \cite{SylvesterUhlmann1987}. In the anisotropic formulation the
unknown is a metric, or equivalently a positive conductivity tensor
density, and boundary-fixing diffeomorphisms give an unavoidable
invariance. Lee and Uhlmann proved a boundary determination result and
uniqueness in the real-analytic setting under geometric and
topological hypotheses \cite{LeeUhlmann1989}.
Lassas and Uhlmann subsequently recovered real-analytic manifolds from
local boundary information \cite{LassasUhlmann2001}.
The Poisson embedding approach of Lassas, Liimatainen, and Salo
\cite{LLS20} expresses the geometry through the values of harmonic
solutions. We use the same broad principle, with source solutions of
a connection Laplacian providing both points and fiber identifications.

For connections, Albin, Guillarmou, Tzou, and Uhlmann recovered
Hermitian connections and matrix potentials from full Cauchy data on
fixed surfaces \cite{AlbinGuillarmouTzouUhlmann2013}. In higher
dimensions, Ceki\'c established recovery for line bundles on certain
known conformally transversally anisotropic geometries
\cite{Cekic2017Connections}, and for smooth unitary Yang--Mills
connections of arbitrary rank from same-patch measurements on an
arbitrary fixed smooth metric \cite[Theorem~1.3]{Cekic2020YM}.
His complex parallel transport method also recovers arbitrary-rank
connections and matrix potentials from full boundary data on fixed
geometries conformal to subdomains of $\R^2\times M_0$
\cite[Theorem~1.1]{Cekic2025Systems}. In the real-analytic category,
Gabdurakhmanov and Kokarev proved joint recovery of metric, Euclidean
bundle, and compatible connection from a boundary patch
\cite[Theorem~1.1]{GabdurakhmanovKokarev2025}. Their use of Green kernels
to identify the base and fibers is a close geometric antecedent.
Theorem~\ref{thm:main} concerns an unknown smooth metric and an arbitrary
smooth unitary connection on the fixed trivial rank-two bundle; it
imposes no Yang--Mills equation on that connection.

The scalar companion \cite{ScalarCompanion2026} develops a smooth
continuation argument based on separately harmonic Green kernels and
shrinking geodesic spheres. The present proof uses its product
continuation, sphere geometry, and scalar angular estimate
\cite[Sections~3--5]{ScalarCompanion2026}. We give the analytic arguments
needed here, including their extension to systems with scalar principal
part. The scalar uniqueness theorem itself is not applied to the
matrix-valued boundary data.

The additional difficulty appears in the first moments of the transfer
between local solution spaces. In scalar harmonic coordinates those
moments can be normalized away. For a connection, a harmonic frame
removes the zeroth-order term, but the transfer still has a
matrix-valued first moment. We represent it by matrices $D^1,\ldots,D^n$.
Their components complementary to the real scalar matrices satisfy a
first-order system whose principal
symbol is the conformal Killing symbol. In dimension at least three
this symbol is of finite type: two prolongations determine all third
derivatives from lower jets. The real scalar part of the displacement
is fixed by the coordinate normalization. Together these facts turn
an initial square-root estimate on $D$ into the stronger estimate
needed to continue the metric and connection across a frontier point.
This reduction is the main matrix-specific step of the proof.

\subsection{How the proof is organized}
Boundary symbol factorization first determines all metric and
connection jets on $\Gamma$, after imposing normal gauge. Attaching a
small exterior cap then converts the boundary form into equality of
Green matrices on an open interior set. Source solutions supported in
that set separate points, span fiber values, and realize the allowed
harmonic jets. They define a maximal local isometry with a unitary
bundle identification; Section~\ref{sec:boundary} prepares coordinates
and harmonic frames at any putative frontier point.

The local continuation mechanism has three stages. First,
Section~\ref{sec:product} continues a mixed Green matrix, initially
known when either variable lies in the matched region. Two
quantitative unique-continuation estimates, one in each variable,
allow continuation across suitable hypersurfaces in the product.
Second, Sections~\ref{sec:spheres} and~\ref{sec:angular} use this kernel
to transfer harmonic functions across small moving spheres. The
transfer is represented by a matrix density of total mass $I$.
An angular coercivity estimate bounds that density uniformly even as
the spheres shrink. The matrix lower-order terms are absorbed in the
same estimate.

Finally, Section~\ref{sec:matrix} extracts the first and second moments
of the transfer. The finite-type argument described above improves
the difference of the normalized inverse metrics from order $r$ to
order $r^{3/2}$, where $r$ is the sphere radius. A continuity argument
then lets the radii collapse, proving local agreement of the
normalized operators. Section~\ref{sec:global} removes the remaining
conformal factor and extends the resulting isometry and unitary
identification over the whole manifold, including every component of
the measured patch.
